\documentclass[10pt,a4paper]{amsart}
\usepackage[margin=19mm]{geometry}
\usepackage{amsmath,amssymb,mathtools}
\usepackage[scr=rsfs,cal=euler]{mathalfa}
\usepackage{xfrac}
\usepackage[unicode,hidelinks,bookmarks=false]{hyperref}
\newcommand{\ie}{i.e.\ }
\newcommand{\cf}{cf.\ }
\newcommand{\resp}{respectively\ }
\newcommand{\Subsection}{\subsection}
\providecommand{\nobreakdash}{\nobreak\hbox{-}}
\setlength{\emergencystretch}{2em}
\multlinegap0pt
\usepackage{amssymb}
\usepackage{dsfont}
\usepackage[normalem]{ulem}
\usepackage{float}
\usepackage{xcolor}
\usepackage{multirow}
\usepackage{enumerate}
\usepackage{caption}
\usepackage[shortlabels]{enumitem}
\usepackage{tikz}
\newtheorem{lemma}{Lemma}
\title{Powers of Hamiltonian cycles in randomly augmented P\'osa-Seymour graphs}
\author{Sylwia Antoniuk, Andrzej Dudek, Andrzej Ruci\'nski}
\date{Source: arXiv:2512.23886v1 (CC BY 4.0)}
\begin{document}
\maketitle

\noindent\emph{Source and licence.} Powers of Hamiltonian cycles in randomly augmented P\'osa-Seymour graphs. Version arXiv:2512.23886v1 (CC BY 4.0), distributed by arXiv under the Creative Commons Attribution 4.0 International licence, http://creativecommons.org/licenses/by/4.0/, as recorded in the arXiv licence metadata for this exact version and shown on the abstract page https://arxiv.org/abs/2512.23886v1. Full licence notice: \texttt{licenses/arxiv-2512.23886-CC-BY-4.0.txt}.\par\smallskip

\noindent\textit{Editorial scope.} Counted unit: the article's lemma weak (source0.tex lines 908--914). The article's complete original proof of it is delivered with it (source0.tex lines 1058--1095, one \texttt{proof} environment of 38 lines, which the article places away from the statement). No mathematics is written, repaired or re-derived: the statement and the proof are the article's own lines, in the article's own notation, and no retained line is substituted or re-flowed.  Retained uncounted as the source-local context the proof needs: lines 396--396; lines 1044--1052.  Nothing was omitted from any retained span, so the package carries no omission record.  No retained line contains a citation, so the wrapper carries no bibliography and no bibliography entry.  No retained line contains a figure environment, an image-inclusion command or drawing code, so no image is delivered and the package declares no asset dependency.  Editorial formatting only, in addition: an AMS \texttt{amsart} preamble that restates the article's own declarations and macros used by the retained lines, namely the environments \texttt{lemma} on separate counters, together with the standard packages those lines need.  The retained environments print in this document as Lemma 1 in order of appearance, whereas the article prints the same items with the numbering of its own full document.  The complete native source of the article as distributed in the arXiv e-print for this exact version is bundled unchanged as \texttt{sources/191837/source0.tex}.
\subsection{New results}\label{newres}

\begin{lemma}\label{perform}
Given an $m$-path $P$ and a subset $V_0\subset V(P)$, algorithm REWIRE outputs a~modified path $P'$ with $V(P')=V(P)$, $|E(P'[V_0])|\le |E(P[V_0])|$, and a $V_0$-based partition $\overrightarrow{V(P')} = T_0\cup S_1\cup T_1\cup\cdots\cup S_q\cup T_q$ such that, setting $x_i=|S_i|$ and $y_i=|T_i|$, $i\in[q]$, we have
\begin{itemize}
    \item[(i)] $1 \leq x_i \leq m$ for $i\in[q]$,
    \item[(ii)] $0 \leq y_i \leq m$ for $i\in[q-1]$,
    \item[(iii)] $x_i + y_i \geq m$ for $i\in[q-1]$,
    \item[(iv)] $y_i + x_{i+1} \geq m$ for $i\in[q-2]$.
\end{itemize}
\end{lemma}

\begin{lemma}\label{weak}
For $k\ge 2$ and $m\geq k+1$ if $P$ is an $m$-path and $V_0\subset V(P)$ satisfies
$|V_0|\ge|V(P)|/(k+1)$, then
\[
    |E(P[V_{0}])|\geq f(\lambda_{k,m})|V_{0}| -  m(m+1).
\]
\end{lemma}

\begin{proof}[Proof of Lemma~\ref{weak}] Let $P$ and $V_0$ be as in the assumptions of the lemma. Setting $L:=|V(P)|$ and $L_0:=|V_0|$, we have
\begin{equation}\label{eq:largest_color}
    k L_0 \ge L - L_0.
\end{equation}

After applying algorithm REWIRE to the pair $(P,V_0)$, we obtain a new path $P'$ with $|E(P[V_0])| \geq |E(P'[V_0])|$ and a $V_0$-bounded partition $\overrightarrow{V(P')}$ satisfying conditions (i)--(iv) of Lemma~\ref{perform}. Condition (iv) does not say anything about $i=q-1$, so we do not know if there is a full $m$-bridge between sets $S_{q-1}$ and $S_q$. To accommodate this deficiency, set $e(S_{q-1},S_q)$ for the number of edges in $P$ with one endpoint in $S_{q-1}$ and the other in $S_q$ and notice that
\begin{equation}\label{lastbridge}
0\le{m+1-y_{q-1}\choose 2}-e(S_{q-1},S_q)\le\binom{m+1}2.
\end{equation}

There is also an issue with condition (ii). Since  we do not have control over how large $y_q$ might be, we subdivide it into terms equal to $m$ plus some remainder, that is, we write
$y_q=y'_q+\cdots+ y'_{q'}$, where $y'_q=\cdots=y'_{q'-1}=m$ and $y'_{q'}\le m$ (when $y_q\le m$, we just have $q'=q$, so nothing changes), and drop the primes over the $y$'s, for convenience. Then ${m+1-y_{i}\choose 2}=0$ for $i=q,\dots, q'-1$, while
\begin{equation}\label{lasty}
{m+1-y_{q'}\choose 2}\le\binom{m+1}2.
\end{equation}
To equalize the number of terms in both sums below, we also, quite artificially, set $x_i:=0$ for all $q+1\le i\le q'$.


 Applying Jensen's inequality to the convex function $\phi(z) = \frac{z(z-1)}{2}$ twice,  with $z_i=x_i$ and $z_i=m+1-y_i$, $i\in[q']$,
 we get
 \begin{equation}\label{Jensen}
 \sum_{i=1}^{q'} {x_i \choose 2}\ge q'{L_0/q' \choose 2}\quad\mbox{and}\quad\sum_{i=1}^{q'} {m+1-y_i\choose 2}\geq q'{m+1-(L-L_0)/q' \choose 2}.
 \end{equation}

 Using (\ref{eq:largest_color})--\eqref{Jensen} and recalling from Section~\ref{newres} the definition of function $f:=f_{k,m}$ with minimum at $\lambda:=\lambda_{k,m}$,
%(taking care of the size of the bridge between $S_{q-1}$ and $S_q$ and the term ${m+1-y_q\choose 2}$ below)
we infer that
\begin{align*}
    |E(P[V_0])| \geq |E(P'[V_0])| & \geq  \sum_{i=1}^q {x_i \choose 2} + \sum_{i=1}^{q-2} {m+1-y_i\choose 2}+e(S_{q-1},S_q) \\
    &\overset{\eqref{lastbridge},\eqref{lasty}}{\geq }\sum_{i=1}^{q'} {x_i \choose 2} + \sum_{i=1}^{q'} {m+1-y_i\choose 2} - m(m+1)\\
    &\hspace{-0.15cm} \overset{\eqref{Jensen}}{\geq}  q'{L_0/q' \choose 2} + q'{m+1-(L-L_0)/q' \choose 2} -  m(m+1) \\
    &\hspace{-0.15cm} \overset{(\ref{eq:largest_color})}{\ge} q'{L_0/q'\choose 2} + q'{m+1-kL_0/q'\choose 2} -  m(m+1) \\
    & =  f(L_0/q')L_0 - m(m+1) \geq f(\lambda)L_0 - m(m+1).
\end{align*}
%where
 %$$c(m) \ge {m+1-y_{q-1}\choose 2}-e(S_{q-1},S_q) +{m+1-y_{q'}\choose 2},$$
  %so, indeed, $c(m)=m(m+1)$ will do.
\end{proof}

\end{document}
